\ifdefined\gkver\else\def\gkver{student}\fi
\documentclass[\gkver]{gaokaozhenti}
\begin{document}

\chapter{2006年全国理综Ⅱ}
%% paperType: 理综物理部分

\begin{enumerate}
\item
%% number: 14
%% paperName: 2006年全国理综Ⅱ第 14 题 6 分
%% typeId: 1
%% type: 单选题
%% chapter: 原子和原子核
%% point: 核聚变
%% method: 无
%% score: 6
%% degree: 650
%% duplicateId: 0
%% body:
现有三个核反应：① \ce{^{24}_{11}Na -> ^{24}_{12}Mg + ^{0}_{-1}e} ② \ce{^{235}_{92}U + ^{1}_{0}n -> ^{141}_{56}Ba + ^{92}_{36}Kr + 3^{1}_{0}n} ③ \ce{^{2}_{1}H + ^{3}_{1}H -> ^{4}_{2}He + ^{1}_{0}n} 下列说法正确的是\xzanswer{C}

\fourchoices[answer=C]
{①是裂变，②是 $\beta$ 衰变，③是聚变}
{①是聚变，②是裂变，③是 $\beta$ 衰变}
{①是 $\beta$ 衰变，②是裂变，③是聚变}
{①是 $\beta$ 衰变，②是聚变，③是裂变}

%% answer: C
%% memo:
\memoanswer{
①核反应放出电子，是 $\beta$ 衰变；②核反应是 \ce{^{235}_{92}U} 在中子的轰击下产生裂变，属于重核裂变；③核反应是氘核和氚核聚变成氦核，属轻核的聚变，故 C 正确。
}

\item
%% number: 15
%% paperName: 2006年全国理综Ⅱ第 15 题 6 分
%% typeId: 1
%% type: 单选题
%% chapter: 力物体的平衡
%% point: 物体系平衡
%% method: 整体与隔离
%% score: 6
%% degree: 650
%% duplicateId: 0
%% body:
如图，位于水平桌面上的物块 P，由跨过定滑轮的轻绳与物块 Q 相连，从滑轮到 P 和到 Q 的两段绳都是水平的。已知 Q 与 P 之间以及 P 与桌面之间的动摩擦因数都是 $\mu$，两物块的质量都是 $m$，滑轮的质量、滑轮轴上的摩擦都不计，若用一水平向右的力 $F$ 拉 P 使它做匀速运动，则 $F$ 的大小为\xzanswer{A}

\onepicture[w=6.0cm]{figs/15.png}

\fourchoices[answer=A]
{$4\mu mg$}
{$3\mu mg$}
{$2\mu mg$}
{$\mu mg$}

%% answer: A
%% memo:
\memoanswer{
P 匀速运动，则 Q 也匀速运动，对 Q 而言，受绳子拉力 $T$ 和 P 对 Q 的摩擦力 $f_{Q}$，绳子的拉力大小等于摩擦力的大小，$T=f_{Q}=\mu mg$；对 P 而言，受向右的拉力 $F$，绳子向左的拉力 $T$，Q 对 P 向左的摩擦力 $f_{Q}$，地面对 P 向左的摩擦力 $f_{P}$，由平衡可知，拉力 $F=T+f_{Q}+f_{P}=\mu mg+\mu mg+\mu\cdot2mg=4\mu mg$，故 A 正确。
}

\item
%% number: 16
%% paperName: 2006年全国理综Ⅱ第 16 题 6 分
%% typeId: 1
%% type: 单选题
%% chapter: 机械波
%% point: 多普勒效应
%% method: 无
%% score: 6
%% degree: 650
%% duplicateId: 0
%% body:
频率一定的声源在空气中向着静止的接收器匀速运动。以 $u$ 表示声源的速度，$V$ 表示声波的速度（$u<V$），$v$ 表示接收器接收到的频率。若 $u$ 增大，则\xzanswer{B}

\fourchoices[answer=B]
{$v$ 增大，$V$ 增大}
{$v$ 增大，$V$ 不变}
{$v$ 不变，$V$ 增大}
{$v$ 减少，$V$ 不变}

%% answer: B
%% memo:
\memoanswer{
当声源向着接收器运动的速度增大时，声源的频率不变，声波的传播速度不变，接收器接收到的频率增大，故 B 正确。
}

\item
%% number: 17
%% paperName: 2006年全国理综Ⅱ第 17 题 6 分
%% typeId: 1
%% type: 单选题
%% chapter: 电场
%% point: 电场叠加
%% method: 对称法
%% score: 6
%% degree: 650
%% duplicateId: 0
%% body:
ab 是长为 $l$ 的均匀带电细杆，P$_{1}$、P$_{2}$ 是位于 ab 所在直线上的两点，位置如图所示。ab 上电荷产生的静电场在 P$_{1}$ 处的场强大小为 $E_{1}$，在 P$_{2}$ 处的场强大小为 $E_{2}$。则以下说法正确的是\xzanswer{D}

\onepicture[w=7.0cm]{\includesvg{figs/17}}

\fourchoices[answer=D]
{两处的电场方向相同，$E_{1}>E_{2}$}
{两处的电场方向相反，$E_{1}>E_{2}$}
{两处的电场方向相同，$E_{1}<E_{2}$}
{两处的电场方向相反，$E_{1}<E_{2}$}

%% answer: D
%% memo:
\memoanswer{
假设杆带正电，可将杆看成由无数个小带电体并排组成，由电场的叠加原理可知，P$_{1}$ 处的电场方向向左，P$_{2}$ 处的电场方向向右，两处的电场方向相反；又由对称性，P$_{1}$ 左侧电荷产生的电场要与右侧抵消一部分，因此 $E_{1}<E_{2}$，故 D 正确。
}

\item
%% number: 18
%% paperName: 2006年全国理综Ⅱ第 18 题 6 分
%% typeId: 1
%% type: 单选题
%% chapter: 运动和力
%% point: 动量守恒定律
%% method: 无
%% score: 6
%% degree: 650
%% duplicateId: 0
%% body:
如图所示，位于光滑水平桌面上的小滑块 P 和 Q 都可视作质点，质量相等。Q 与轻质弹簧相连。设 Q 静止，P 以某一初速度向 Q 运动并与弹簧发生碰撞。在整个碰撞过程中，弹簧具有的最大弹性势能等于\xzanswer{B}

\onepicture[w=6.0cm]{figs/18.png}

\fourchoices[answer=B]
{P 的初动能}
{P 的初动能的 $\frac{1}{2}$}
{P 的初动能的 $\frac{1}{3}$}
{P 的初动能的 $\frac{1}{4}$}

%% answer: B
%% memo:
\memoanswer{
由题意，当两滑块具有相同速度时，弹簧被压缩至最短，弹簧具有的弹性势能最大，系统动能的减小量等于弹簧弹性势能的增加量，由动量守恒 $mv_{0}=2mv$，得 $v=\frac{v_{0}}{2}$，由能量守恒得 $E_{p}=\frac{1}{2}mv_{0}^{2}-\frac{1}{2}\cdot2m\left(\frac{v_{0}}{2}\right)^{2}=\frac{1}{4}mv_{0}^{2}=\frac{1}{2}E_{k0}$，因此最大弹性势能等于 P 初动能的 $\frac{1}{2}$，故 B 正确。
}

\item
%% number: 19
%% paperName: 2006年全国理综Ⅱ第 19 题 6 分
%% typeId: 2
%% type: 多选题
%% chapter: 光学
%% point: 光电效应
%% method: 无
%% score: 6
%% degree: 650
%% duplicateId: 0
%% body:
已知能使某金属产生光电效应的极限频率为 $\nu_{0}$，则\xzanswer{AB}

\fourchoices[answer=AB]
{当用频率为 $2\nu_{0}$ 的单色光照射该金属时，一定能产生光电子}
{当用频率为 $2\nu_{0}$ 的单色光照射该金属时，所产生的光电子的最大初动能为 $h\nu_{0}$}
{当照射光的频率 $\nu$ 大于 $\nu_{0}$ 时，若 $\nu$ 增大，则逸出功增大}
{当照射光的频率 $\nu$ 大于 $\nu_{0}$ 时，若 $\nu$ 增大一倍，则光电子的最大初动能也增大一倍}

%% answer: AB
%% memo:
\memoanswer{
某金属产生光电效应的极限频率为 $\nu_{0}$，则逸出功 $W=h\nu_{0}$；当用频率为 $2\nu_{0}$ 的单色光照射该金属时，一定能产生光电子，故 A 正确；由光电效应方程得 $E_{k}=2h\nu_{0}-W$，有 $E_{k}=h\nu_{0}$，故 B 正确；逸出功只与金属的材料有关，与照射光的频率无关，故 C 错误；若照射光的频率增大一倍，逸出功并没有增大一倍，光电子的最大初动能增加大于一倍，故 D 错误。
}

\item
%% number: 20
%% paperName: 2006年全国理综Ⅱ第 20 题 6 分
%% typeId: 2
%% type: 多选题
%% chapter: 电磁感应
%% point: 电磁感应能量分析
%% method: 无
%% score: 6
%% degree: 650
%% duplicateId: 0
%% body:
如图所示，位于一水平面内的、两根平行的光滑金属导轨，处在匀强磁场中，磁场方向垂直于导轨所在的平面，导轨的一端与一电阻相连；具有一定质量的金属杆 ab 放在导轨上并与导轨垂直。现用一平行于导轨的恒力 $F$ 拉杆 ab，使它由静止开始向右运动。杆和导轨的电阻、感应电流产生的磁场均可不计。用 $E$ 表示回路中的感应电动势，$i$ 表示回路中的感应电流，在 $i$ 随时间增大的过程中，电阻消耗的功率等于\xzanswer{BD}

\onepicture[w=6.0cm]{\includesvg{figs/20}}

\fourchoices[answer=BD]
{$F$ 的功率}
{安培力的功率的绝对值}
{$F$ 与安培力的合力的功率}
{$iE$}

%% answer: BD
%% memo:
\memoanswer{
$F$ 为恒力，杆做变速运动，先加速，后匀速，感应电流先逐渐增大，由功能关系，外力做的功等于杆动能的增加和电阻上消耗的电能，电阻上消耗的电能等于杆克服安培力做的功，因此电阻上消耗的功率等于安培力功率的绝对值，故 B 正确；又 $P=UI$，在 $i$ 随时间增大的过程中，电阻上消耗的功率为瞬时功率，所以 $P=Ei$，故 D 正确。
}

\item
%% number: 21
%% paperName: 2006年全国理综Ⅱ第 21 题 6 分
%% typeId: 1
%% type: 单选题
%% chapter: 气体性质
%% point: 盖吕萨克定律
%% method: 无
%% score: 6
%% degree: 650
%% duplicateId: 0
%% body:
对一定量的气体，若用 $N$ 表示单位时间内与器壁单位面积碰撞的分子数，则\xzanswer{C}

\fourchoices[answer=C]
{当体积减小时，$N$ 必定增加}
{当温度升高时，$N$ 必定增加}
{当压强不变而体积和温度变化时，$N$ 必定变化}
{当压强不变而体积和温度变化时，$N$ 可能不变}

%% answer: C
%% memo:
\memoanswer{
单位时间内与器壁单位面积碰撞的分子数与气体的密集程度和分子的平均动能有关，即与气体的体积和温度有关，当体积减小时，单位体积内的分子数增加，但温度变化未知，$N$ 不一定增加，故 A 错误；当温度升高时，分子的平均动能增大，但气体体积可能增大，$N$ 不一定增加，故 B 错误；当压强不变时，即单位时间内器壁单位面积受到的压力不变，体积和温度变化时，分子平均动能变化，而压强不变，则 $N$ 必定变化，故 C 正确，D 错误。
}

\item
%% number: 22
%% paperName: 2006年全国理综Ⅱ第 22 题 11 分
%% typeId: 4
%% type: 实验题
%% chapter: 光学
%% point: 测定玻璃折射率
%% method: 无
%% score: 11
%% degree: 650
%% duplicateId: 0
%% body:
\begin{enumerate}
	\item
	（6 分）现要测定一个额定电压 $4\UV$、额定功率 $1.6\UW$ 的小灯泡（图中用 $\otimes$ 表示）的伏安特性曲线。要求所测电压范围为 $0.1\UV\sim4\UV$。

	现有器材：直流电源 E（电动势 $4.5\UV$，内阻不计），电压表 V（量程 $4.5\UV$，内阻约为 $4\times10^{4}\UO$），电流表 A$_{1}$（量程 $250\UmA$，内阻约为 $2\UO$），电流表 A$_{2}$（量程 $500\UmA$，内阻约为 $1\UO$），滑动变阻器 $R$（最大阻值约为 $30\UO$），电子键 S，导线若干。

	如果既要满足测量要求，又要测量误差较小，应该选用的电流表是 \tkanswer[1.5cm]{A$_{2}$}，下面两个电路应该选用的是 \tkanswer[1.5cm]{甲}。

	\twopicture[numstyle=chinese, labelprefix={},
		labelA=2006全国理综Ⅱ22a, labelB=2006全国理综Ⅱ22b, wA=2.8cm, wB=2.8cm]
	{figs/22a.png}
	{figs/22b.png}

	\item
	（11 分）一块玻璃砖用两个相互平行的表面，其中一个表面是镀银的（光线不能通过表面）。现要测定此玻璃的折射率。给定的器材还有：白纸、铅笔、大头针 4 枚（P$_{1}$、P$_{2}$、P$_{3}$、P$_{4}$）、带有刻度的直角三角板、量角器。

	实验时，先将玻璃砖放到白纸上，使上述两个相互平行的表面与纸面垂直。在纸上画出直线 aa$'$ 和 bb$'$，aa$'$ 表示镀银的玻璃表面，bb$'$ 表示另一表面，如图所示。然后，在白纸上竖直插上两枚大头针 P$_{1}$、P$_{2}$（位置如图）。用 P$_{1}$、P$_{2}$ 的连线表示入射光线 。

	i. 为了测量折射率，应如何正确使用大头针 P$_{3}$、P$_{4}$？\tkanswer[6cm]{}。

	试在题图中标出 P$_{3}$、P$_{4}$ 的位置。

	ii. 然后，移去玻璃砖与大头针。试在题图中通过作图的方法标出光线从空气到玻璃中的入射角 $\theta_{1}$ 与折射角 $\theta_{2}$。简要写出作图步骤。\tkanswer[6cm]{}。

	iii. 写出 $\theta_{1}$、$\theta_{2}$ 表示的折射率公式为 $n=$ \tkanswer[2cm]{}。

	\onepicture[h=5.5cm, align=right]{figs/22c.png}
\end{enumerate}

%% answer:
\jdanswer{
\begin{enumerate}
	\item
	A$_{2}$，甲。
	\item
	i. 在 bb$'$ 一侧观察 P$_{1}$、P$_{2}$（经 bb$'$ 折射、aa$'$ 反射，再经 bb$'$ 折射后）的像，在适当的位置插上 P$_{3}$，使得 P$_{3}$ 与 P$_{1}$、P$_{2}$ 的像在一条直线上，即让 P$_{3}$ 挡住 P$_{1}$、P$_{2}$ 的像；再插上 P$_{4}$，让它挡住 P$_{2}$（或 P$_{1}$）的像和 P$_{3}$。P$_{3}$、P$_{4}$ 的位置如图所示。

	\onepicture[w=6.0cm]{figs/22_an.png}

	ii. ①过 P$_{1}$、P$_{2}$ 作直线与 bb$'$ 交于 O；②过 P$_{3}$、P$_{4}$ 作直线与 bb$'$ 交于 O$'$；③利用刻度尺找到 OO$'$ 的中点 M；④过 O 点作 bb$'$ 的垂线 CD，过 M 点作 bb$'$ 的垂线与 aa$'$ 相交于 N 点，如图所示，连接 ON；⑤$\angle$P$_{1}$OD $=\theta_{1}$，$\angle$CON $=\theta_{2}$。

	iii. $n=\frac{\sin\theta_{1}}{\sin\theta_{2}}$
\end{enumerate}
}

%% memo:
\memoanswer{
（1）小灯泡的额定功率 $1.6\UW$，当电压调为最大等于额定电压 $4\UV$ 时，流过小灯泡的电流为 $400\UmA$，故电流表应选 A$_{2}$；又因小灯泡的电阻为 $10\UO$，比较小，接近电流表 A$_{2}$ 的内阻，电流表应采用外接法，故选择甲图。

（2）i. 在 bb$'$ 一侧观察 P$_{1}$、P$_{2}$（经 bb$'$ 折射、aa$'$ 反射，再经 bb$'$ 折射后）的像，在适当的位置插上 P$_{3}$，使得 P$_{3}$ 与 P$_{1}$、P$_{2}$ 的像在一条直线上，即让 P$_{3}$ 挡住 P$_{1}$、P$_{2}$ 的像；再插上 P$_{4}$，让它挡住 P$_{2}$（或 P$_{1}$）的像和 P$_{3}$，P$_{3}$、P$_{4}$ 的位置如图所示。

ii. ①过 P$_{1}$、P$_{2}$ 作直线与 bb$'$ 交于 O；②过 P$_{3}$、P$_{4}$ 作直线与 bb$'$ 交于 O$'$；③利用刻度尺找到 OO$'$ 的中点 M；④过 O 点作 bb$'$ 的垂线 CD，过 M 点作 bb$'$ 的垂线与 aa$'$ 相交于 N，如图所示，连接 ON；⑤$\angle$P$_{1}$OD $=\theta_{1}$，$\angle$CON $=\theta_{2}$。

iii. $\frac{\sin\theta_{1}}{\sin\theta_{2}}$
}

\item
%% number: 23
%% paperName: 2006年全国理综Ⅱ第 23 题 16 分
%% typeId: 6
%% type: 计算题
%% chapter: 曲线运动
%% point: 平抛运动
%% method: 无
%% score: 16
%% degree: 650
%% duplicateId: 0
%% body:
如图所示，一固定在竖直平面内的光滑的半圆形轨道 ABC，其半径 $R=0.5\Um$，轨道在 C 处与水平地面相切。在 C 处放一小物块，给它一水平向左的初速度 $v_{0}=5\Ums$，结果它沿 CBA 运动，通过 A 点，最后落在水平面上的 D 点，求 C、D 间的距离 $s$。

\onepicture[h=4.2cm, align=right]{\includesvg{figs/23}}

%% answer:
\jdanswer{$s=1\Um$}

%% memo:
\memoanswer{
设小物块的质量为 $m$，过 A 处时的速度为 $v$，由 A 到 D 经历的时间为 $t$，由动能定理得

$\frac{1}{2}mv_{0}^{2}=\frac{1}{2}mv^{2}+2mgR$ ①

由平抛运动的规律得

$2R=\frac{1}{2}gt^{2}$ ②

$s=vt$ ③

由①②③式并代入数据得

$s=1\Um$ ④
}

\item
%% number: 24
%% paperName: 2006年全国理综Ⅱ第 24 题 19 分
%% typeId: 6
%% type: 计算题
%% chapter: 运动和力
%% point: 超重失重
%% method: 无
%% score: 19
%% degree: 650
%% duplicateId: 0
%% body:
一质量为 $m=40\Ukg$ 的小孩子站在电梯内的体重计上。电梯从 $t=0$ 时刻由静止开始上升，在 0 到 $6\Us$ 内体重计示数 $F$ 的变化如图所示。试问：在这段时间内电梯上升的高度是多少？

\onepicture[h=4.5cm, align=right]{\includesvg{figs/24}}

%% answer:
\jdanswer{$h=9\Um$}

%% memo:
\memoanswer{
由图可知，在 $t=0$ 到 $t=t_{1}=2\Us$ 的时间内，体重计的示数大于 $mg$，故电梯应做向上的加速运动。设这段时间内体重计作用于小孩的力为 $F_{1}$，电梯及小孩的加速度为 $a_{1}$，由牛顿第二定律，得

$F_{1}-mg=ma_{1}$ ①

在这段时间内电梯上升的高度

$h_{1}=\frac{1}{2}a_{1}t_{1}^{2}$ ②

在 $t_{1}$ 到 $t=t_{2}=5\Us$ 的时间内，体重计的示数等于 $mg$，故电梯应做匀速上升运动，速度为 $t_{1}$ 时刻电梯的速度，即

$v_{1}=a_{1}t_{1}$ ③

在这段时间内电梯上升的高度

$h_{2}=v_{1}(t_{2}-t_{1})$ ④

在 $t_{2}$ 到 $t=t_{3}=6\Us$ 的时间内，体重计的示数小于 $mg$，故电梯应做向上的减速运动。设这段时间内体重计作用于小孩的力为 $F_{2}$，电梯及小孩的加速度为 $a_{2}$，由牛顿第二定律，得

$mg-F_{2}=ma_{2}$ ⑤

在这段时间内电梯上升的高度

$h_{3}=v_{1}(t_{3}-t_{2})-\frac{1}{2}a_{2}(t_{3}-t_{2})^{2}$ ⑥

电梯上升的总高度

$h=h_{1}+h_{2}+h_{3}$ ⑦

由以上各式，利用牛顿第三定律和题文及题图中的数据，解得

$h=9\Um$ ⑧
}

\item
%% number: 25
%% paperName: 2006年全国理综Ⅱ第 25 题 20 分
%% typeId: 6
%% type: 计算题
%% chapter: 磁场
%% point: 带电粒子在磁场中的运动
%% method: 无
%% score: 20
%% degree: 450
%% duplicateId: 0
%% body:
如图所示，在 $x<0$ 与 $x>0$ 的区域中，存在磁感应强度大小分别为 $B_{1}$ 与 $B_{2}$ 的匀强磁场，磁场方向均垂直于纸面向里，且 $B_{1}>B_{2}$。一个带负电荷的粒子从坐标原点 O 以速度 $v$ 沿 $x$ 轴负方向射出，要使该粒子经过一段时间后又经过 O 点，$B_{1}$ 与 $B_{2}$ 的比值应满足什么条件？

\onepicture[w=4.6cm, align=right]{figs/25.png}

%% answer:
\jdanswer{$\frac{B_{2}}{B_{1}}=\frac{n}{n+1}$（$n=1,2,3,\ldots$）}

%% memo:
\memoanswer{
粒子在整个运动过程中的速度大小恒为 $v$，交替地在 $xy$ 平面内 $B_{1}$ 与 $B_{2}$ 磁场区域中做匀速圆周运动，轨道都是半个圆周。设粒子的质量和电荷量的大小分别为 $m$ 和 $q$，圆周运动的半径分别为 $r_{1}$ 和 $r_{2}$，有

$r_{1}=\frac{mv}{qB_{1}}$ ①

$r_{2}=\frac{mv}{qB_{2}}$ ②

现分析粒子运动的轨迹。如图所示，在 $xy$ 平面内，粒子先沿半径为 $r_{1}$ 的半圆 $C_{1}$ 运动至 $y$ 轴上离 O 点距离为 $2r_{1}$ 的 A 点，接着沿半径为 $r_{2}$ 的半圆 $D_{1}$ 运动至 O$_{1}$ 点，OO$_{1}$ 的距离

$d=2(r_{2}-r_{1})$ ③

此后，粒子每经历一次“回旋”（即从 $y$ 轴出发沿半径为 $r_{1}$ 的半圆和半径为 $r_{2}$ 的半圆回到原点下方的 $y$ 轴），粒子的 $y$ 坐标就减小 $d$。设粒子经过 $n$ 次回旋后与 $y$ 轴交于 O$_{n}$ 点，若 OO$_{n}$ 即 $nd$ 满足

$nd=2r_{1}$ ④

则粒子再经过半圆 $C_{n+1}$ 就能经过原点，式中 $n=1,2,3,\cdots$ 为回旋次数。

由③④式解得

$\frac{r_{1}}{r_{2}}=\frac{n}{n+1}$（$n=1,2,3,\cdots$） ⑤

联立①②⑤式可得 $B_{1}$、$B_{2}$ 应满足的条件：

$\frac{B_{2}}{B_{1}}=\frac{n}{n+1}$（$n=1,2,3,\cdots$） ⑥

\onepicture[w=6.0cm]{figs/25_an.png}
}

\end{enumerate}

\end{document}
